% TOP_PROBLEM: {"id": 30007198, "title": "Coefficient-preserving solutions of nonsingular group-equation systems", "category_id": 17, "category": {"id": 17, "name": "group_theory", "display_name": "Group Theory", "description": "Problems about groups, group actions, representations, and related algebraic structures.", "slug": "group-theory", "order_index": 17, "created_at": "2026-07-31T15:26:25.671Z"}, "status": "open"}
% FIRST_POSTED: 2026-10-06
% ATTEMPT_STATUS: solved
% OPENAI_CLAIM: {"schema_version":1,"source_repo":"https://github.com/openai/math","source_commit":"adc7f1241b42e322a6451854ab7e4b4c146bf78a","release_date":"2026-10-06","match":"full","resolution":"proved","independently_reviewed":false,"second_pass":"agreed","families":[{"family":"256","title":"Nonsingular systems of equations over arbitrary groups","subject":"Group theory","manuscripts":[{"title":"Nonsingular systems of equations over arbitrary groups","pdf_path":"preprints/Nonsingular-systems-of-equations-over-arbitrary-groups-October-5-2026/nonsingular-systems-over-arbitrary-groups.pdf","readme_path":"preprints/Nonsingular-systems-of-equations-over-arbitrary-groups-October-5-2026/README.md","date":"2026-10-05","theorem_ref":"thm:main (sections/introduction.tex:27)","statement_sources":[{"path":"preprints/Nonsingular-systems-of-equations-over-arbitrary-groups-October-5-2026/build/sections/introduction.tex","line":27,"label":"thm:main"}]}],"lean":{"doc_path":"lean/docs/256.md","comparators":[],"declarations":[],"covers_main_theorem":false},"reasoning_trace":null}]}
\documentclass[11pt]{article}
\usepackage[margin=25mm]{geometry}
\usepackage{fontspec}
\setmainfont{Latin Modern Roman}
\usepackage{amsmath,amssymb,mathtools}
\usepackage{unicode-math}
\setmathfont{Latin Modern Math}
\usepackage{enumitem}
\usepackage{xurl,hyperref,bookmark}
\hypersetup{colorlinks=true,urlcolor=blue}
\setcounter{secnumdepth}{0}
\setlength{\parindent}{0pt}
\setlength{\parskip}{5pt}
\setlength{\emergencystretch}{3em}
\begin{document}
\section*{Coefficient-preserving solutions of nonsingular group-equation systems}

\subsection{OpenAI's claim}
OpenAI presents a proof in the selected manuscript. Any finite equation system over a group with full-rank exponent-sum matrix is claimed simultaneously solvable after embedding into a larger group.

\subsection{Mathematical statement of the claimed result}
\paragraph{Result 1.} Let \(G\) be any group, \(1\le m\le n\), and \(w_i\in G*F_n\). For the exponent-sum matrix \(E=(e_{ij})\), where \(e_{ij}\) counts the total exponent of the \(j\)-th free generator in \(w_i\), assume \(\operatorname{rank}_{\mathbb Q}E=m\). The source claims injectivity of
\[
 G\longrightarrow (G*F_n)/\langle\!\langle w_1,\ldots,w_m\rangle\!\rangle.
\]
Equivalently, all equations \(w_i=1\) can be solved simultaneously in an overgroup without identifying distinct coefficient elements of \(G\).

\subsection{Outline of the argument}
Simultaneous unitary fixed-space constructions are measured by a traced spectral phase. Subadditivity and weighted-cycle calculations rule out the planar coefficient obstruction associated with a nonsingular system, allowing the coefficient group to embed in the solution overgroup.

\subsection{Scope relative to this problem}
The October theorem has arbitrary G, 1<=m<=n and rational row rank m, and explicitly proves coefficient injectivity and equivalent simultaneous overgroup solvability. lean/docs/256.md only links the earlier one-variable unimodular paper, not the systems theorem. Within these assumptions, this is a claimed affirmative answer to the stated question.

\subsection{Formal verification}
For family 256, the selected statement is coefficient injectivity for a single relator of exponent sum $1$ or $-1$ in $A*\mathbb Z$. It underlies the Kervaire consequence but does not include the newer general nonsingular-system theorem. No Comparator is assigned to the selected manuscript here; the family page selects other papers or supporting statements. This site has not executed the Lean code. The scope descriptions do not certify the remaining manuscript arguments.

\subsection{Status caveat}
OpenAI attributes these manuscripts to work by an internal model. Its release cautions that arguments without formalization can contain mistakes. This catalogue has not independently checked the proofs; the labels report the mathematical claims attributed to OpenAI.

\subsection{Sources}
\begin{itemize}
\item OpenAI, \emph{Nonsingular systems of equations over arbitrary groups}, 2026-10-05; citation key \texttt{OAI:Nonsingular-systems-of-equations-over-arbitrary-groups-October-5-2026}. \href{https://github.com/openai/math/blob/adc7f1241b42e322a6451854ab7e4b4c146bf78a/preprints/Nonsingular-systems-of-equations-over-arbitrary-groups-October-5-2026/nonsingular-systems-over-arbitrary-groups.pdf}{PDF (GitHub)}; \href{https://github.com/openai/math/raw/adc7f1241b42e322a6451854ab7e4b4c146bf78a/preprints/Nonsingular-systems-of-equations-over-arbitrary-groups-October-5-2026/nonsingular-systems-over-arbitrary-groups.pdf}{PDF (direct)}; \href{https://github.com/openai/math/blob/adc7f1241b42e322a6451854ab7e4b4c146bf78a/preprints/Nonsingular-systems-of-equations-over-arbitrary-groups-October-5-2026/README.md}{Manuscript page and citation}; \href{https://github.com/openai/math/blob/main/preprints/Nonsingular-systems-of-equations-over-arbitrary-groups-October-5-2026/nonsingular-systems-over-arbitrary-groups.pdf}{Latest version}; Statement source: \href{https://github.com/openai/math/blob/adc7f1241b42e322a6451854ab7e4b4c146bf78a/preprints/Nonsingular-systems-of-equations-over-arbitrary-groups-October-5-2026/build/sections/introduction.tex\#L27}{thm:main}.
\item Family 256: \href{https://github.com/openai/math/blob/adc7f1241b42e322a6451854ab7e4b4c146bf78a/lean/docs/256.md}{Lean scope}.
\item \href{https://github.com/openai/math/blob/adc7f1241b42e322a6451854ab7e4b4c146bf78a/CONTENTS.md}{OpenAI catalogue entry}.
\end{itemize}
\end{document}
